\subsection{A family with the path structure}\label{sec:mechanism}

The benchmark models do not tell whether the cuts can supply strength that
SCIP lacks where the theory predicts a gap. We therefore built a family with
the structure of Section~\ref{sec:interleave}. An instance with $n$ copies has
variables $x_i,y_i,z_i\in[0,1]$ and $t_i\in[-10,10]$, the rows
$\Phi^{(i)}(x_i,y_i,z_i)-t_i\le0$, where $\Phi^{(i)}$ is the
function~\eqref{eq:family} with $\omega_A=\omega_C=1$, one coupling row
$\sum_iy_i\le0.8n$, and the objective $\min\sum_it_i$. For each copy, four
distinct values are drawn from $\{k/64:k=0,\dots,48\}$ and split into two
interleaving two-point sets $A_i$ and $C_i$ with random orientation. All
coefficients are dyadic and exactly representable. Every optimal $y_i$ is a
midpoint of two values of at most $3/4$, so the coupling row does not bind, and
by Theorem~\ref{thm:interleave} the optimum is $\sum_i\delta_i^2/2$. Glued exact
pair hulls give the bound $0$ for every copy: the chords of $A_i$ and $C_i$
cross at a point with $m_{y_i}\le\max A_i\le3/4$, so the zero points of all
copies together satisfy the coupling row. The part of the gap that only the
joint block closes is therefore the whole optimum. Part~3C used five instances
for each $n\in\{10,20,40,80\}$ (seeds 0--4); Part~4C3 adds five fresh instances
per $n$ (seeds 5--9), which no earlier run had used. The randomized family was
adopted after a pilot with SCIP alone had shown that the deterministic family,
with every copy equal to the example of Proposition~\ref{prop:path-gap}, is
easy for SCIP: it solved $160$ copies in 19\,s.

Because the objective is a sum of block functions and the coupling row does not
bind, the $n$ cuts $t_i\ge\min\Phi^{(i)}$, which are the exact supports of the
blocks in their whole-row directions, make the root bound equal to the
optimum. This is the most favorable case for block cuts. Part~4C4 removes it:
it replaces the coupling row of each fresh instance by $\sum_iy_i\le c$, where
$c$ is half the sum of the smallest unconstrained minimizers $y_i^\ast$, rounded
down to a multiple of $1/64$. The row then binds, the block minima no longer
give the optimum, and the cuts must be found in other directions. For Part~4C4
we know two reference values to high accuracy. The first is the optimum, from
a convex mixed-integer reformulation (one binary variable per choice of nearest
points) solved by Gurobi, re-solved exactly for the returned choice and
certified by an exact branch and bound to within $10^{-58}$. The second is the
\emph{block closure}: the bound of all aggregated cuts on the blocks
$(x_i,y_i,z_i)$ together with the linear rows. Each block has the single
nonlinear row $\Phi^{(i)}-t_i\le0$ and the domain $[0,1]^3$, because the coupling
row involves other blocks and, since $c\ge1$, implies no tighter bounds. By
Theorem~\ref{thm:closure} and Corollary~\ref{cor:one-row} the closure of the
block is $t_i\ge\operatorname{vex}\Phi^{(i)}(x_i,y_i,z_i)$; minimizing over $x_i$
and $z_i$ commutes with convexification, because the projection of a convex
hull is the convex hull of the projection. With
$d_i(y)=\min_{x,z\in[0,1]}\Phi^{(i)}(x,y,z)=\dist(y,A_i)^2+\dist(y,C_i)^2$ the
block closure is therefore
$\min\{\sum_i\operatorname{vex}d_i(y_i):\sum_iy_i\le c,\ 0\le y\le1\}$, which is
also the Lagrangian dual bound of the coupling row. We computed it with the
exact convex envelopes of the piecewise quadratic functions $d_i$ and a
bisection on the multiplier, to within $2\cdot10^{-28}$.
Table~\ref{tab:mechanism} summarizes the results; Tables~\ref{tab:path-detail}
and~\ref{tab:path-fresh} in Appendix~\ref{app:instances} give them per
instance.

\begin{table}[t]
\centering\footnotesize\setlength{\tabcolsep}{4pt}
\caption{The path family: median fraction of SCIP's root gap closed,
$(z_{\text{mode}}-z_{\text{base}})/(z^\ast-z_{\text{base}})$, for each $n$ (five
instances each), and instances solved within 300\,s (of 20). ``Remainder'' and
``whole row'' are the first directions of the separator, ``base'' and ``wide''
its limits (Section~\ref{sec:blocks}, Table~\ref{tab:modes}). ``Glued pair
hulls'' is the analytic bound~0 of Theorem~\ref{thm:interleave} and ``block
closure'' the bound of all cuts on the blocks; no run computes them. SCIP
closes 0 by definition, and Gurobi had no root runs. Rows of different parts
ran under different host loads, so their solved counts are not paired
(Section~\ref{sec:setup}). The binding-row block shows three decimals; its
Part~5C-b rows are root runs only.}\label{tab:mechanism}
\begin{tabular}{@{}lllrrrrr@{}}
\toprule
& & & \multicolumn{4}{c}{Root gap closed (median)} & \\
\cmidrule(lr){4-7}
Instances & Part & Mode & $n=10$ & $n=20$ & $n=40$ & $n=80$ & Solved\\
\midrule
seeds 0--4 & 3C & SCIP & & & & & 9\\
 & 4C2 & SCIP-nolocks & 0.95 & 0.89 & 0.93 & 0.91 & 10\\
 & 4C2 & SCIP-extra & 0.64 & 0.58 & 0.66 & 0.59 & 10\\
 & 4C2 & Gurobi & & & & & 7\\
 & --- & glued pair hulls & 0.95 & 0.89 & 0.93 & 0.92 & \\
 & 3C & remainder, base & 0.35 & 0.25 & 0.19 & 0.22 & 9\\
 & 4C2 & remainder, wide & 0.86 & 0.87 & 0.89 & 0.90 & 13\\
 & 3P & whole row, base & 0.44 & 0.44 & 0.43 & 0.46 & 10\\
 & 3P & whole row, wide & 1.00 & 1.00 & 1.00 & 1.00 & 20\\
\midrule
seeds 5--9 & 4C3 & SCIP & & & & & 9\\
 & 5C-a & SCIP-nolocks & 0.92 & 0.92 & 0.92 & 0.92 & 10\\
 & 4C3 & SCIP-extra & 0.57 & 0.64 & 0.61 & 0.66 & 9\\
 & 4C3 & Gurobi & & & & & 5\\
 & --- & glued pair hulls & 0.93 & 0.92 & 0.93 & 0.92 & \\
 & 4C3 & remainder, base & 0.32 & 0.31 & 0.24 & 0.20 & 10\\
 & 4C3 & whole row, wide & 1.00 & 1.00 & 1.00 & 1.00 & 20\\
\midrule
binding row & 4C4 & SCIP & & & & & 10\\
(seeds 5--9) & 5C-a & SCIP-nolocks & 0.980 & 0.980 & 0.972 & 0.959 & 20\\
 & 4C4 & SCIP-extra & 0.287 & 0.600 & 0.685 & 0.655 & 13\\
 & 4C4 & Gurobi & & & & & 6\\
 & --- & block closure & 1.000 & 1.000 & 1.000 & 1.000 & \\
 & 4C4 & remainder, wide & 0.983 & 0.960 & 0.954 & 0.933 & 20\\
 & 4C4 & whole row, wide & 0.997 & 0.989 & 0.981 & 0.982 & 20\\
 & 5C-b & remainder, cap $64n$ & 1.000 & 0.999 & 1.000 & 0.999 & \\
 & 5C-b & whole row, cap $64n$ & 1.000 & 0.999 & 1.000 & 0.999 & \\
\bottomrule
\end{tabular}
\end{table}

\paragraph{SCIP's root bound and the pair-hull level.} SCIP's root bound is
weak: on the instances of Part~3C it lies between $-0.035$ and $-0.013$ per
copy, while each copy contributes at least $1/8192$ to the optimum. The cause
is a presolve step: $\Phi^{(i)}$ is concave in $x_i$ and in $z_i$, so SCIP's
implicit-discreteness detection (Section~\ref{sec:related}) makes these
variables binary, and SCIP then closes the gap by branching. SCIP-nolocks,
with this step switched off, raises the root bound to just below~0, the bound
of glued exact pair hulls, on every instance of Part~3C; SCIP's semidefinite
minor cuts were productive in all of these runs and in none of the default
runs of campaign~4. On the fresh instances of Part~4C3, SCIP-nolocks
(Part~5C-a) again stopped just below the pair-hull level, with a median
closure of 0.92. SCIP-extra, of whose extra separators only the
intersection cuts were productive, closed about 60\% of SCIP's root gap. None
of the three SCIP settings we ran closed any part of the gap between the
pair-hull level and the optimum, which is the gap that
Section~\ref{sec:composition} attributes to joint blocks. Gurobi~13 solved the
instances of Part~4C2 with $n=10$ and two with $n=20$; on $n=40$ and $n=80$ its
gap after 300\,s was 61 to 113\% (53 to 121\% on the instances of Part~4C3).

\paragraph{The prospective separator and the attribution.} The separator of
Part~3C (remainder directions, base limits) improved the root bound on all 20
instances but closed only a median of 19--35\% of the root gap, less than the
pair-hull level. It solved the same 9 instances as SCIP, with fewer nodes and a
shifted geometric mean time of 6.7\,s instead of 10.3\,s. After seeing these
results we identified two limitations. First, the remainder directions bound
only the nonlinear remainder of a row, whereas the support of the whole row is
the block minimum. Second, the first blocks of a callback take up to four cuts
each, so with $n$ cuts per callback the later blocks receive none. A post hoc
diagnostic (Part~3P) tried whole-row directions first, once with base limits
and once with wide limits; we had tried the wide limits by hand on two of the
20 instances before. With base limits, whole-row directions raised the median
closure to 43--46\% and solved 10 instances; with wide limits they closed the
root gap on all 20 instances and solved all 20, 18 of them at the root node.
Part~4C2 ran the remaining combination prospectively: remainder directions
with wide limits closed 86--90\%, close to the pair-hull level and above it on
7 of the 20 instances, and solved 13. Every run of these four configurations
stopped at its cut cap, so their root bounds do not depend on the host load,
which was lower in campaign~4. The wide limits therefore account for most of
the closure; on this family the whole-row directions supply the rest and raise
the solved count from 13 to 20.

\paragraph{Prospective test on fresh instances.} On the 20 instances of
Part~4C3, the whole-row variant with wide limits, fixed before these instances
were generated, closed the root gap on every instance: its root bound was
within $2\cdot10^{-7}$ of the optimum. It solved all 20 in 3 to 33\,s, with a
median of one node and at most 38, against 9 for SCIP, 9 for SCIP-extra,
10 for SCIP-nolocks, 10 for the separator with remainder directions and base
limits, and 5 for Gurobi. On the nine instances that SCIP solved, the
cut mode was slower on four of the five with $n=10$, by factors of up to 3.4
because of the separator's fixed cost, and 4 to 20 times faster on the four
with $n=20$; SCIP's own time fell to a median of 4\% of its time in the
baseline. Over all 20 runs, 96\% of the time was spent in the Python separator.
The separator with remainder directions and base limits closed 20--32\% of the
root gap, as in Part~3C.

\paragraph{A binding coupling row (Part 4C4).} When the coupling row binds,
the block minima no longer give the optimum, and the cuts must have the right
slopes in $y_i$. Both separator configurations with wide limits solved all 20
instances, in 3 to 134\,s, against 10 for SCIP, 13 for SCIP-extra and 6 for
Gurobi; on the ten instances that SCIP solved, all with $n\le20$, the cut modes
were slower, by median factors of 3.5 (remainder directions) and 3.2
(whole-row directions), because of the separator's time. At the root the two
cut modes closed medians of 95\% (remainder directions) and 99\% (whole-row
directions) of SCIP's root gap. The block closure equals the optimum on 13 of
the 20 instances and lies at most $6.5\cdot10^{-4}$ below it on the others. No
root run reached it: every cut-mode run stopped at its cap of $16n$ cuts, and
with whole-row directions the root bound stayed between 0.016 and 0.029 below
the block closure for $n=80$. With the cap $64n$ and up to 40 callbacks
(Part~5C-b), separation ended by itself after 14 to 27 callbacks, with medians
of 29 (remainder directions) and 26 (whole-row directions) cuts per block, and
the root bound came within $8\cdot10^{-4}$ of the block closure on every
instance; with the cap $32n$ the runs were the same, except that the
remainder-direction runs reached the cap on three instances. The distance left
with the $16n$ cap was therefore due to the cap, not to the direction search.
In the runs with the $16n$ cap, 99.9\% (remainder directions) and 94\%
(whole-row directions) of the cuts came from the LP direction search, almost
all with a nonzero coefficient of $y_i$, mostly with the negative slope that the
binding row requires; none of them is an envelope cut in $(y_i,t_i)$ alone. The
whole-row cuts alone give only $\sum_i\min d_i$, which is below SCIP's root bound
on half of the instances.

SCIP-nolocks (Part~5C-a), however, also solved all 20 instances, in 0.1 to
169\,s. Its shifted geometric mean time was 2.0\,s, against 12 to 13\,s for the
cut modes in Part~4C4; these times come from different campaigns and are not
paired, but the difference is much larger than the effect of the host load.
Its root bound closed a median of 96 to 98\% of SCIP's root gap. With a binding
row, most of the gap is thus within reach of SCIP's own relaxation once
implicit discreteness is switched off, and the part that only the joint blocks
close is small.

Both path families have a convex mixed-integer reformulation, the one we used
to compute the optima, which Gurobi solves in under a second on 19 of the 20
instances of Part~4C4; the experiment therefore tests whether a solver
recovers the block structure from the nonconvex formulation. With a single
coupling row, the Shapley--Folkman estimate of \citet{AubinEkeland1976} bounds
the distance between the block closure and the optimum by twice the largest
nonconvexity $\max_y(d_i(y)-\operatorname{vex}d_i(y))$ of a single block,
independently of $n$; the observed distances were much smaller still. These
families were constructed to have the structure of
Section~\ref{sec:composition}, and their objectives are sums of block
functions, which favors block cuts. On the family with a non-binding row,
exact block support supplied root strength that none of the three SCIP
settings obtained, and the separator solved instances that neither SCIP nor
Gurobi solved within 300\,s; this strength reached the solver only when the
separator could add enough cuts per block in useful directions. With a binding
row, SCIP without implicit discreteness was as successful and faster. The
families are not evidence about other models.
